\documentclass[12pt,a4paper]{article}

\usepackage[utf8]{inputenc}
\usepackage[T1]{fontenc}
\usepackage[spanish,es-tabla]{babel}
\usepackage{amsmath,amssymb}
\usepackage{geometry}
\usepackage{pgfplots}
\pgfplotsset{compat=1.18}

\geometry{margin=2.5cm}

\title{Antitransformada de Laplace y representación temporal}
\author{}
\date{}

\begin{document}

\maketitle

\section*{Función dada}

Se considera la función de transferencia o transformada:

\begin{equation}
    Y(s)=\frac{s-1}{(s+2)(s+4)}.
\end{equation}

Se busca obtener la señal $y(t)$ en el dominio del tiempo y representarla gráficamente.

\section*{Descomposición en fracciones simples}

Planteamos:

\begin{equation}
    \frac{s-1}{(s+2)(s+4)}
    =\frac{A}{s+2}+\frac{B}{s+4}.
\end{equation}

Multiplicando por $(s+2)(s+4)$:

\begin{align}
    s-1 &= A(s+4)+B(s+2),\\
    s-1 &= (A+B)s+(4A+2B).
\end{align}

Igualando coeficientes:

\begin{equation}
    \begin{cases}
        A+B=1,\\
        4A+2B=-1.
    \end{cases}
\end{equation}

De donde:

\begin{equation}
    A=-\frac{3}{2},\qquad B=\frac{5}{2}.
\end{equation}

Por lo tanto:

\begin{equation}
    Y(s)=-\frac{3}{2}\frac{1}{s+2}
    +\frac{5}{2}\frac{1}{s+4}.
\end{equation}

\section*{Expresión en el dominio del tiempo}

Usando la pareja de transformadas de Laplace

\begin{equation}
    \mathcal{L}^{-1}\left\{\frac{1}{s+a}\right\}=e^{-at}u(t),
\end{equation}

se obtiene:

\begin{equation}
    \boxed{
    y(t)=\left(-\frac{3}{2}e^{-2t}+\frac{5}{2}e^{-4t}\right)u(t)
    }.
\end{equation}

Para $t\geq 0$:

\begin{equation}
    \boxed{
    y(t)=-\frac{3}{2}e^{-2t}+\frac{5}{2}e^{-4t}
    }.
\end{equation}

\section*{Características de la señal}

La señal comienza en:

\begin{equation}
    y(0^+)=1.
\end{equation}

El cruce por cero ocurre cuando:

\begin{align}
    -\frac{3}{2}e^{-2t}+\frac{5}{2}e^{-4t}&=0,\\
    t&=\frac{1}{2}\ln\left(\frac{5}{3}\right)\approx 0.2554.
\end{align}

Además,

\begin{equation}
    \lim_{t\to\infty}y(t)=0,
\end{equation}

aproximándose al cero desde valores negativos.

\section*{Gráfica}

\begin{center}
\begin{tikzpicture}
\begin{axis}[
    width=14cm,
    height=8cm,
    axis lines=middle,
    xlabel={$t$},
    ylabel={$y(t)$},
    xmin=0, xmax=3,
    ymin=-0.16, ymax=1.08,
    xtick={0,0.5,1,1.5,2,2.5,3},
    ytick={-0.1,0,0.2,0.4,0.6,0.8,1},
    grid=both,
    minor tick num=1,
    samples=200,
    domain=0:3,
    smooth,
    thick,
    blue
]
    \addplot {-(3/2)*exp(-2*x)+(5/2)*exp(-4*x)};
    \addplot[only marks, mark=*, mark size=1.6pt, red]
        coordinates {(0,1) (0.2554,0)};
\end{axis}
\end{tikzpicture}
\end{center}

\section*{Resultado final}

\begin{equation}
    \boxed{
    Y(s)=\frac{s-1}{(s+2)(s+4)}
    \quad\Longleftrightarrow\quad
    y(t)=\left(-\frac{3}{2}e^{-2t}+\frac{5}{2}e^{-4t}\right)u(t)
    }.
\end{equation}

\end{document}
